\section{Rangkuman}
\begin{itemize}
    \item Alice dan Bob perlu memegang kunci AES yang sama agar dapat bertukar
          pesan terenkripsi. Ada dua jalan keluar: kriptografi hibrida atau
          algoritma pertukaran kunci Diffie--Hellman.
    \item Kriptografi hibrida menggabungkan kriptografi kunci-simetri dengan
          kriptografi kunci-publik. Alice mengenkripsi kunci rahasia $K$ dengan
          RSA memakai kunci publik Bob,
          $C_K = E_{\text{RSA},\text{PK}_{\text{Bob}}}(K)$, lalu mengenkripsi
          pesan dengan AES, $C_M = E_{\text{AES},K}(M)$.
    \item Bob memulihkan kunci dengan kunci privatnya,
          $K = D_{\text{RSA},\text{SK}_{\text{Bob}}}(C_K)$, lalu memulihkan pesan
          dengan AES, $M = D_{\text{AES},K}(C_M)$.
    \item Algoritma pertukaran kunci Diffie--Hellman dipublikasikan oleh Diffie
          dan Hellman pada tahun 1976. Protokol ini dipakai untuk berbagi kunci
          rahasia yang sama, lalu kunci itu dipakai oleh algoritma
          kunci-simetri.
    \item Parameter publiknya adalah bilangan prima $p$ dan bilangan bulat $g$
          dengan $g < p$. Bilangan $g$ harus akar primitif dari $p$. Nilai $p$
          dan $g$ tidak perlu rahasia.
    \item Prosedurnya: Alice mengirim $A = g^a \bmod p$, Bob mengirim
          $B = g^b \bmod p$, lalu keduanya menghitung
          $K = B^a = A^b = g^{ab} \bmod p$.
    \item Contoh dengan $p = 353$, $g = 3$, $a = 97$, dan $b = 233$ menghasilkan
          $A = 40$, $B = 248$, dan $K = 160$. Contoh perkakas daring dengan
          $G = 1601$ dan $N = 4789$ menghasilkan $A = 2716$, $B = 1302$, dan
          $K = 4257$.
    \item Keamanan Diffie--Hellman bersandar pada sulitnya menghitung logaritma
          diskret. Penyadap yang mengetahui $p$, $g$, $A$, dan $B$ tetap sukar
          menemukan $a$ atau $b$ bila $p$ sangat besar.
    \item Diffie--Hellman aman terhadap penyadapan, tetapi tidak tahan terhadap
          serangan aktif \textit{man-in-the-middle}. Mallory mengganti $g^a$ dan
          $g^b$ dengan $g^m$, sehingga ia memegang $K_1 = g^{am} \bmod p$ dan
          $K_2 = g^{bm} \bmod p$.
    \item Protokol ini dapat diperluas ke tiga pihak dengan pola melingkar.
          Alice, Bob, dan Carol sampai pada kunci yang sama,
          $K = g^{abc} \bmod p$.
    \item Diffie--Hellman dipakai di dalam protokol SSL dan TLS. Pada tahun
          2002, Hellman menyarankan nama Diffie--Hellman--Merkle untuk menghargai
          kontribusi Ralph Merkle.
\end{itemize}
